% Cut block: Autocovariances and spectrum (Section 4)
\subsection{Autocovariances and spectrum}
\label{subsec:spectrum}

The same transition matrix gives the complete two-point statistics by finite-dimensional spectral decomposition.

\begin{theorem}\label{thm:spectrum}
The autocovariances \(c_j:=\Cov(\Delta_0,\Delta_j)\) satisfy
\[
c_0=\frac{20}{81},
\]
and, for \(j\ge 1\),
\[
c_j=2^{-j}\!\left(\frac{1}{8}+(-1)^j\!\left(\frac{7}{648}+\frac{j}{108}\right)\right).
\]
The two-sided power spectral density
\[
S(\omega):=\sum_{j\in \ZZ} c_j\,e^{-i\omega j}
\]
is the rational function
\[
S(\omega)
=
\frac{2\bigl(32\cos^3\omega+12\cos^2\omega-116\cos\omega-105\bigr)}
{9\,(4\cos\omega-5)\,(4\cos\omega+5)^2}.
\]
In particular,
\[
S(0)=\sum_{j\in\ZZ} c_j=\frac{118}{243},
\]
in agreement with Corollary~\ref{thm:llnclt}.
\end{theorem}

\begin{proof}
Appendix~\ref{app:pressure} expresses the covariances as
\[
c_j=qP^{j-1}h-\left(\frac49\right)^2\qquad (j\ge1),
\]
for explicit vectors \(q\) and \(h\) attached to the Parry transition matrix \(P\). Because \(P\) is similar to \(A/2\), its nontrivial spectral modes are \(1/2\) and \(-1/2\), with a double root at \(-1/2\). Hence \(c_j\) is a linear combination of \(2^{-j}\), \((-2)^{-j}\), and \(j(-2)^{-j}\), and matching the first three values yields the displayed formula. The series is absolutely summable, so the power spectral density is obtained by summing the resulting cosine series. Appendix~\ref{app:pressure} performs that sum explicitly and obtains the stated rational function.
\end{proof}


% Cut block: Joint anomaly--output rotation set (Section 4)
\subsection{The joint anomaly--output rotation set}
\label{subsec:joint-rotation}

The preceding marginal laws do not describe which output densities and
discrepancy densities can occur simultaneously. As a further consequence of
the pair presentation, the standard cycle method makes this constraint exact.
For a shift-invariant Borel probability measure
\(\mu\) on \(\mathcal P\), set
\[
a(\mu):=\int Y_0\,d\mu,
\qquad
d(\mu):=\int (X_0\oplus Y_0)\,d\mu.
\]

\begin{theorem}[Sharp anomaly--output tradeoff]
\label{thm:joint-rotation}
The joint rotation set is the pentagon
\begin{align}
\mathcal R
&:=\left\{(a(\mu),d(\mu)):
\text{\(\mu\) is shift-invariant on \(\mathcal P\)}\right\}\notag\\
&=\operatorname{conv}\left\{
(0,0),\ \left(\frac12,0\right),
\left(\frac12,\frac12\right),
\left(\frac13,1\right),
\left(0,\frac12\right)
\right\}\label{eq:joint-rotation-polygon}\\
&=\left\{(a,d)\in\RR^2:
0\le a\le\frac12,\quad 0\le d,\quad
d\le\frac12+\frac32a,\quad d\le2-3a
\right\}.\notag
\end{align}
In particular, the largest discrepancy density compatible with output density
\(a\in[0,1/2]\) is
\begin{equation}
d_{\max}(a)=
\begin{cases}
\dfrac12+\dfrac32a,&0\le a\le\dfrac13,\\[2mm]
2-3a,&\dfrac13\le a\le\dfrac12.
\end{cases}
\label{eq:max-discrepancy-profile}
\end{equation}

The measures attaining the first upper bound in
\eqref{eq:joint-rotation-polygon} are exactly those supported on the subshift
\(\mathcal P_{\mathrm L}\) presented by
\[
00\xrightarrow{01}01,\qquad
01\xrightarrow{10}10,\qquad
10\xrightarrow{10}00,\qquad
10\xrightarrow{00}01.
\]
The measures attaining the second upper bound are exactly those supported on
the subshift \(\mathcal P_{\mathrm R}\) obtained by replacing the last edge by
\(10\xrightarrow{11}01\). Both frontier subshifts are mixing and satisfy
\[
h_{\mathrm{top}}(\mathcal P_{\mathrm L})
=h_{\mathrm{top}}(\mathcal P_{\mathrm R})
=\log\rho_{\mathrm{pl}},
\qquad
\rho_{\mathrm{pl}}^3-\rho_{\mathrm{pl}}-1=0,
\]
where \(\rho_{\mathrm{pl}}>1\) is the plastic constant. Each vertex of
\(\mathcal R\) is realized by a unique invariant measure, supported on the
corresponding periodic pair word
\[
(00)^\ZZ,\quad (11|00)^\ZZ,\quad (10|11)^\ZZ,
\quad (01|10|10)^\ZZ,\quad (10|00)^\ZZ,
\]
listed in the order of the vertices in \eqref{eq:joint-rotation-polygon}.
\end{theorem}

\begin{proof}
Attach to every edge \(e\) of \(G_{\mathrm{pair}}\) the vector
\[
v(e):=(Y(e),X(e)\oplus Y(e)).
\]
The edge shift of \(G_{\mathrm{pair}}\) maps onto \(\mathcal P\) by a
finite-to-one factor map. Every invariant measure on \(\mathcal P\) has an
invariant lift to the edge shift, and every invariant edge measure projects
to \(\mathcal P\). The one-edge marginal of a lifted measure is a stationary
flow on the finite graph. Every such flow is a convex combination of
normalized directed-cycle flows; this is the cycle description of rotation
sets for locally constant observables \cite{Ziemian1995}. Conversely, each
directed cycle supports an invariant periodic-orbit measure. It is therefore
enough to list the simple directed cycles, retaining the two choices on the
parallel edges from state \(10\) to state \(01\):
\[
\begin{array}{c|c|c}
\text{cycle labels}&a&d\\ \hline
00&0&0\\
11|00&1/2&0\\
10|11&1/2&1/2\\
01|10|10&1/3&1\\
10|00&0&1/2
\end{array}
\]
These are all simple cycles of the four-state graph. Their rotation vectors
are precisely the five points in \eqref{eq:joint-rotation-polygon}; taking
their convex hull gives the first description of \(\mathcal R\). Computing
the five supporting lines gives the inequality description and
\eqref{eq:max-discrepancy-profile}.

For the supporting functional \(d-3a/2\), the maximizing simple cycles are
exactly \(10|00\) and \(01|10|10\). The union of their edges is the graph
presenting \(\mathcal P_{\mathrm L}\). A cycle decomposition shows that a
stationary edge flow attains the supporting value \(1/2\) if and only if it
assigns full mass to this union. The same argument for the supporting
functional \(d+3a\) shows that its maximizing cycles are exactly
\(10|11\) and \(01|10|10\), whose edge union presents
\(\mathcal P_{\mathrm R}\). This proves both equality classifications.

In the state order \((00,01,10)\), either frontier graph has adjacency matrix
\[
C=
\begin{pmatrix}
0&1&0\\
0&0&1\\
1&1&0
\end{pmatrix}.
\]
It is primitive, and
\[
\det(\lambda I-C)=\lambda^3-\lambda-1.
\]
The right-resolving presentations are finite-to-one on label words, so their
label entropy equals the edge entropy \(\log\rho(C)=\log\rho_{\mathrm{pl}}\).
Finally, each vertex in the cycle table is represented by only one simple
cycle. Equality in a convex decomposition at an extreme point therefore
forces the corresponding periodic-orbit measure, proving uniqueness.
\end{proof}


% Cut block: Weighted-fiber thermodynamics (Section 6)
\section{Finite-graph consequences for weighted fibers}
\label{sec:partition}

The four-state presentation from Section~\ref{sec:gauge} also keeps track of the fibers of \(\Fold_m\). No further structural input is needed: weighting its edges gives the recurrence, output-weight asymptotics, and variational interpretation below. The linear-algebraic eliminations are placed in Appendix~\ref{app:recurrence}.

\subsection{Weighted fibers}

For \(x\in X_m\), define the fiber multiplicity
\[
d_m(x):=\#\Fold_m^{-1}(x).
\]
The corresponding weighted partition function is
\begin{equation}\label{eq:Zm}
Z_m(y):=\sum_{x\in X_m} y^{\abs{x}_1}\, d_m(x).
\end{equation}
Equivalently,
\[
Z_m(y)=\sum_{\omega\in\Om_m} y^{\abs{\Fold_m(\omega)}_1},
\]
so \(Z_m(y)\) is the moment-generating polynomial of the output Hamming weight under the uniform measure on \(\Om_m\). We also write
\[
Z(t;y):=\sum_{m\ge 0} Z_m(y)\,t^m,
\qquad
Z_0(y)=1.
\]

\subsection{Transfer matrix}

Reading the four-state graph of Appendix~\ref{app:pressure} as a weighted automaton, one assigns to an edge labeled by \((X,Y)\in\{0,1\}^2\) the weight \(y^Y\). This records the contribution of that transition to \(\abs{\Fold_m(\omega)}_1\). Choosing state \(00\) as the initial carry state gives a four-dimensional linear representation; the terminal vector below incorporates the finite-window boundary conditions:
\begin{equation}\label{eq:boundary-vector}
Z_m(y)=\alpha^{\mathsf T} M(y)^m \beta(y),
\end{equation}
where
\begin{equation}\label{eq:boundary-data}
\alpha=
\begin{pmatrix}
1\\ 0\\ 0\\ 0
\end{pmatrix},
\qquad
\beta(y)=
\begin{pmatrix}
1\\[1mm]
\dfrac{1}{1+y}\\[2mm]
\dfrac{1}{y}\\[2mm]
\dfrac{y}{1+y}
\end{pmatrix},
\end{equation}
and
\begin{equation}\label{eq:My}
M(y)=
\begin{pmatrix}
1&y&0&y\\
0&0&1&0\\
1&1+y&0&0\\
1&0&0&0
\end{pmatrix}.
\end{equation}
Appendix~\ref{app:recurrence} derives \(\beta(y)\) from the initial values \(Z_0,\dots,Z_3\). Although \(\beta(y)\) is rational, the product \(\alpha^{\mathsf T}M(y)^m\beta(y)\) is a polynomial in \(y\); the apparent poles at \(y=0\) and \(y=-1\) cancel identically.

\subsection{Consequences of the transfer matrix}

The following formulas are direct finite-dimensional consequences of the weighted adjacency matrix.

\begin{proposition}\label{prop:recurrence}
The sequence \((Z_m(y))_{m\ge 0}\) satisfies
\begin{equation}\label{eq:Zm-rec}
Z_m(y)=Z_{m-1}(y)+(2y+1)Z_{m-2}(y)-Z_{m-3}(y)-y(y+1)Z_{m-4}(y)
\qquad (m\ge 4),
\end{equation}
with initial values
\[
Z_0(y)=1,\qquad
Z_1(y)=1+y,\qquad
Z_2(y)=2+2y,\qquad
Z_3(y)=2+5y+y^2.
\]
Equivalently,
\begin{equation}\label{eq:Zty}
Z(t;y)=\frac{N(t;y)}{D(t;y)},
\end{equation}
where
\begin{equation}\label{eq:Dty}
D(t;y)=1-t-(2y+1)t^2+t^3+y(y+1)t^4
\end{equation}
and
\begin{equation}\label{eq:Nty}
N(t;y)=(1+ty)(1-yt^2).
\end{equation}
The characteristic polynomial of \(M(y)\) is
\begin{equation}\label{eq:Pi}
\Pi(\lambda,y)=\lambda^4-\lambda^3-(2y+1)\lambda^2+\lambda+y(y+1).
\end{equation}
\end{proposition}

\begin{proof}
Appendix~\ref{app:recurrence} derives \eqref{eq:boundary-vector}. A direct determinant computation gives
\[
\det(\lambda I-M(y))=\Pi(\lambda,y),
\]
hence \eqref{eq:Zm-rec} by Cayley--Hamilton. Summing the recurrence against \(t^m\) and using the initial values yields \eqref{eq:Zty}--\eqref{eq:Nty}.
\end{proof}

Appendix~\ref{app:recurrence} also records the exceptional parameter values at which the rational generating function is not minimal, together with the explicit discriminant factorization of \(\Pi(\lambda,y)\). Those formulas are auxiliary for our purposes. For the asymptotic results below we only use that, for \(y>0\), Perron--Frobenius theory gives a real, positive, simple dominant root \(\lambda_\star(y)\) of \(\Pi(\lambda,y)\). Consequently the free energy
\[
f(y):=\lim_{m\to\infty}\frac1m\log Z_m(y)=\log \lambda_\star(y)
\]
is real-analytic on \((0,\infty)\).

\subsection{Output-weight statistics}
\label{subsec:weight-stats}

Let \(u_{\Om_m}\) be the uniform measure on \(\Om_m\), and let \(w_m:=\Fold_m\push u_{\Om_m}\) be its pushforward to \(X_m\). For \(X\sim w_m\), write \(H_m:=\abs{X}_1\).

\begin{corollary}\label{thm:weight-stats}
Under \(w_m\),
\[
\frac{H_m}{m}\xrightarrow{\ \PP\ }\frac{5}{18},
\qquad
\frac{H_m-\frac{5}{18}m}{\sqrt m}\Rightarrow \mathcal{N}\!\left(0,\frac{67}{972}\right).
\]
Moreover, \((H_m/m)_{m\ge1}\) satisfies a large deviation principle on \([0,1]\) with good rate function
\[
I(a):=\sup_{\theta\in\RR}\bigl(\theta a-\Lambda_Y(\theta)\bigr),
\qquad
\Lambda_Y(\theta):=\log \rho\!\bigl(M(e^\theta)\bigr)-\log 2,
\]
where \(\rho(\cdot)\) denotes the Perron root.
\end{corollary}

\begin{proof}
Since \(w_m=\Fold_m\push u_{\Om_m}\),
\[
\EE_{w_m}[e^{\theta H_m}]
=
2^{-m}\sum_{\omega\in\Om_m} e^{\theta\abs{\Fold_m(\omega)}_1}
=
2^{-m}Z_m(e^\theta).
\]
By Proposition~\ref{prop:recurrence}, \(Z_m(e^\theta)\) is governed by the transfer matrix \(M(e^\theta)\), so Perron--Frobenius theory gives
\[
\lim_{m\to\infty}\frac{1}{m}\log \EE_{w_m}[e^{\theta H_m}]
=
\log \rho\!\bigl(M(e^\theta)\bigr)-\log 2
=
\Lambda_Y(\theta).
\]
The function \(\Lambda_Y\) is real-analytic near \(0\). Differentiating the Perron root at \(\theta=0\) in Appendix~\ref{app:recurrence} gives
\[
\Lambda_Y'(0)=\frac{5}{18},
\qquad
\Lambda_Y''(0)=\frac{67}{972}.
\]
The convergence in probability and the central limit theorem then follow from the standard perturbative theory for primitive matrix families \cite[Chapter~4]{ParryPollicott1990}, \cite[Chapter~6]{Ruelle2004}. Since \(\theta\mapsto M(e^\theta)\) is analytic for all \(\theta\in\RR\), the G\"artner--Ellis theorem yields the stated large deviation principle \cite[Theorem~2.3.6]{DemboZeitouni1998}.
\end{proof}

\subsection{A microcanonical entropy profile}
\label{subsec:microcanonical}

Let \(\mathcal{P}\subset(\{00,01,10,11\})^{\ZZ}\) be the pair shift from Appendix~\ref{app:pressure}. For \(p=(p_t)_{t\in\ZZ}\in\mathcal{P}\), let
\[
\upsilon(p):=\text{the second coordinate of }p_0.
\]
Thus \(\upsilon\) is the locally constant observable whose Birkhoff average is the output density. For \(a\in[0,\frac12]\), define
\[
\mathfrak h(a):=
\sup\left\{
h_\mu(\shift):
\begin{array}{l}
\mu \text{ is a }\shift\text{-invariant Borel probability measure on }\mathcal P,\\[1mm]
\int \upsilon\,d\mu=a
\end{array}
\right\}.
\]

\begin{theorem}\label{thm:microcanonical-output}
The attainable set of output densities is the full interval \([0,\frac12]\). For every \(\theta\in\RR\),
\begin{equation}\label{eq:microcanonical-legendre}
\log \rho\!\bigl(M(e^\theta)\bigr)
=
\sup_{0\le a\le 1/2}\bigl(\mathfrak h(a)+\theta a\bigr).
\end{equation}
The function \(a\mapsto \mathfrak h(a)\) is real-analytic and strictly concave on \((0,\frac12)\), and its unique maximizer is
\[
\mathfrak h\!\left(\frac{5}{18}\right)=\log 2.
\]
Moreover, the rate function from Corollary~\ref{thm:weight-stats} satisfies
\[
I(a)=\log 2-\mathfrak h(a)\qquad\left(0\le a\le \frac12\right),
\]
and
\[
I(a)=+\infty\qquad\left(\frac12<a\le 1\right).
\]
\end{theorem}

\begin{proof}
Every output word belongs to the golden-mean language, so every invariant measure on \(\mathcal P\) satisfies \(\int \upsilon\,d\mu\le 1/2\). The fixed point carried by the loop
\[
00\xrightarrow{00}00
\]
has output density \(0\), while the \(2\)-cycle
\[
00\xrightarrow{11}11\xrightarrow{00}00
\]
has output density \(1/2\). Convex combinations of these two invariant measures therefore realize every density in \([0,\frac12]\).

Appendix~\ref{app:pressure} gives a synchronized right-resolving presentation of \(\mathcal P\), and \(\upsilon\) is a locally constant potential on this sofic shift. Weighting each edge by \(e^{\theta\upsilon}\) yields exactly the matrix \(M(e^\theta)\), so the topological pressure of \(\theta\upsilon\) is
\[
P(\theta):=P_{\mathcal P}(\theta\upsilon)=\log \rho\!\bigl(M(e^\theta)\bigr);
\]
see \cite[Chapter~4]{ParryPollicott1990}, \cite[Chapter~9]{Walters1982}, or \cite[Chapter~6]{Ruelle2004}. By the variational principle,
\[
P(\theta)
=
\sup_{\mu}\left(h_\mu(\shift)+\theta\int \upsilon\,d\mu\right),
\]
where the supremum is over all \(\shift\)-invariant Borel probability measures on \(\mathcal P\). Grouping the measures by the value of \(\int \upsilon\,d\mu\) gives \eqref{eq:microcanonical-legendre}.

The function \(P\) is real-analytic because \(M(e^\theta)\) is an analytic primitive matrix family. It is strictly convex, since \(\upsilon\) is not cohomologous to a constant: its average is \(0\) on the fixed point above and \(1/2\) on the displayed \(2\)-cycle. Standard equilibrium-state theory for mixing sofic shifts therefore gives \(P''(\theta)>0\) for every \(\theta\in\RR\); see \cite[Chapter~4]{ParryPollicott1990} or \cite[Chapter~9]{Walters1982}.

For \(\theta\le 0\), one has \(0\le \upsilon\le 1\), hence \(0\le P(\theta)\le P(0)=\log 2\). For \(\theta\ge 0\), every length-\(n\) pair block is determined by its \(n\) microstate digits, so there are at most \(2^n\) such blocks, and the corresponding output has at most \(\lceil n/2\rceil\) ones. Hence the partition sum defining pressure is at most \(2^n e^{\theta\lceil n/2\rceil}\), which yields
\[
P(\theta)\le \log 2+\frac{\theta}{2}.
\]
The \(2\)-cycle above gives the lower bound \(P(\theta)\ge \theta/2\) for \(\theta\ge0\). Therefore
\[
\lim_{\theta\to-\infty}P'(\theta)=0,
\qquad
\lim_{\theta\to+\infty}P'(\theta)=\frac12.
\]
Legendre--Fenchel duality then shows that \(\mathfrak h\) is real-analytic and strictly concave on \((0,\frac12)\), with
\[
\mathfrak h(a)=\inf_{\theta\in\RR}\bigl(P(\theta)-\theta a\bigr).
\]
At \(\theta=0\), the equilibrium state is the Parry measure of the pair graph, so
\[
P(0)=\log 2,
\qquad
P'(0)=\Lambda_Y'(0)=\frac{5}{18}
\]
by Appendix~\ref{app:pressure}. Hence the unique maximizer of \(\mathfrak h\) is \(a=5/18\), and its value is \(\log 2\).

Finally,
\[
I(a)=\sup_{\theta\in\RR}\bigl(\theta a-\Lambda_Y(\theta)\bigr)
=
\sup_{\theta\in\RR}\bigl(\theta a-P(\theta)+\log 2\bigr).
\]
For \(0\le a\le 1/2\), the dual formula for \(\mathfrak h\) yields \(I(a)=\log 2-\mathfrak h(a)\). If \(a>1/2\), then every \(x\in X_m\) satisfies \(\abs{x}_1\le m/2\), so the probabilities in Corollary~\ref{thm:weight-stats} vanish identically on a neighborhood of \(a\). Thus \(I(a)=+\infty\) on \((1/2,1]\).
\end{proof}

Theorem~\ref{thm:microcanonical-output} identifies the weighted theory with an entropy maximization problem on the pair shift. In this form, the transfer matrix governing \(Z_m(y)\) and the variational description of the free energy are simply two descriptions of the same finite-state model.

For comparison, interval-based Zeckendorf statistics \cite{KologluKoppMillerWang2011,MillerWang2012,BowerInsoftLiMillerTosteson2015} study the decomposition of a uniformly chosen integer in \([1,\Fib_n]\). In that model the summand density converges to \(1/(\varphi^2+1)\approx 0.2764\), whereas the fold-pushforward measure \(w_m\) yields the density \(5/18\approx 0.2778\). The closeness of these constants appears to be accidental, since the underlying sampling measures are different.


% Cut block: Tilted matrices, cumulants, autocovariances, and power spectrum (Appendix B)
\subsection{Tilted matrices and cumulants}
For the discrepancy observable \(\Delta_t=X_t\oplus Y_t\), only the labels \(01\) and \(10\) contribute. The tilted adjacency matrix is therefore
\[
A_\theta=
\begin{pmatrix}
1& e^\theta &0&1\\
0&0& e^\theta &0\\
e^\theta &2&0&0\\
1&0&0&0
\end{pmatrix}.
\]
Write \(z=e^\theta\), and let \(\rho_\Delta(\theta)\) be the Perron root. Then
\[
F_\Delta(\lambda,z)
:=
\det(\lambda I-A_\theta)
=
\lambda^4-\lambda^3-(2z+1)\lambda^2+(2z-z^3)\lambda+2z.
\]
Since \(F_\Delta(2,1)=0\), implicit differentiation of
\[
F_\Delta(\rho_\Delta(\theta),e^\theta)=0
\]
at \(\theta=0\) gives
\[
F_{\Delta,\lambda}(2,1)=9,\qquad
F_{\Delta,z}(2,1)=-8,
\]
so \(\rho_\Delta'(0)=8/9\). A second differentiation yields
\[
F_{\Delta,\lambda\lambda}(2,1)=30,\qquad
F_{\Delta,\lambda z}(2,1)=-9,\qquad
F_{\Delta,zz}(2,1)=-12,
\]
hence \(\rho_\Delta''(0)=332/243\). Therefore
\[
\Lambda_\Delta(\theta):=\log \rho_\Delta(\theta)-\log 2
\]
satisfies
\[
\Lambda_\Delta'(0)=\frac49,
\qquad
\Lambda_\Delta''(0)=\frac{118}{243}.
\]
These are the mean discrepancy density and asymptotic variance used in Corollary~\ref{thm:llnclt}.

For the output-weight observable \(Y_t\), the labels with \(Y_t=1\) are \(01\) and \(11\). Weighting each such edge by \(e^\theta\) gives
\[
M(e^\theta)=
\begin{pmatrix}
1& e^\theta &0& e^\theta\\
0&0&1&0\\
1&1+e^\theta&0&0\\
1&0&0&0
\end{pmatrix},
\]
the weighted matrix used in Section~\ref{sec:partition}. Its characteristic polynomial is
\[
F_Y(\lambda,z)
:=
\det(\lambda I-M(z))
=
\lambda^4-\lambda^3-(2z+1)\lambda^2+\lambda+z(z+1).
\]
Since \(F_Y(2,1)=0\), implicit differentiation of \(F_Y(\rho_Y(\theta),e^\theta)=0\) gives
\[
F_{Y,\lambda}(2,1)=9,\qquad
F_{Y,z}(2,1)=-5,
\]
so \(\rho_Y'(0)=5/9\). Differentiating again and using
\[
F_{Y,\lambda\lambda}(2,1)=30,\qquad
F_{Y,\lambda z}(2,1)=-8,\qquad
F_{Y,zz}(2,1)=2,
\]
gives \(\rho_Y''(0)=71/243\). Hence
\[
\Lambda_Y(\theta):=\log \rho_Y(\theta)-\log 2
\]
satisfies
\[
\Lambda_Y'(0)=\frac{5}{18},
\qquad
\Lambda_Y''(0)=\frac{67}{972}.
\]

\subsection{Autocovariances and power spectrum}
For the vertex chain with transition matrix \(P\), the discrepancy observable depends only on the current transition:
\[
g(i,j)=
\begin{cases}
1,& (i,j)\in\{(00,01),(01,10),(10,00)\},\\
0,& \text{otherwise}.
\end{cases}
\]
Define
\[
h_j:=\sum_k P_{jk}g(j,k)
=
\begin{pmatrix}
\frac14\\[2mm]
1\\[2mm]
\frac12\\[2mm]
0
\end{pmatrix},
\qquad
q_k:=\sum_j \pi_jP_{jk}g(j,k)
=
\begin{pmatrix}
\frac19 & \frac19 & \frac29 & 0
\end{pmatrix}.
\]
Then
\[
\EE[\Delta_0]=\pi h=\frac49,
\qquad
\EE[\Delta_0\Delta_j]=qP^{j-1}h\qquad (j\ge1).
\]
Because \(P=\tfrac12 D^{-1}AD\) with \(D=\mathrm{diag}(2,1,2,1)\), \(P\) is similar to \(A/2\). Therefore its characteristic polynomial is
\[
(\lambda-1)\left(\lambda-\frac12\right)\left(\lambda+\frac12\right)^2.
\]
The stationary covariance sequence \(c_j=\Cov(\Delta_0,\Delta_j)\) therefore has the form
\[
c_j=2^{-j}\bigl(a+(-1)^j(b+cj)\bigr)\qquad (j\ge1).
\]
Direct evaluation of \(qP^{j-1}h-(4/9)^2\) at \(j=1,2,3\) gives
\[
c_1=\frac{17}{324},\qquad
c_2=\frac{25}{648},\qquad
c_3=\frac{7}{648},
\]
from which
\[
a=\frac18,\qquad
b=\frac{7}{648},\qquad
c=\frac{1}{108}.
\]
Hence
\[
c_j = 2^{-j}\!\left(\frac{1}{8}+(-1)^j\!\left(\frac{7}{648}+\frac{j}{108}\right)\right)\qquad (j\ge1),
\]
while
\[
c_0=\Var(\Delta_0)=\frac49\left(1-\frac49\right)=\frac{20}{81}.
\]
Substituting these values into
\[
S(\omega)=c_0+2\sum_{j\ge1}c_j\cos(j\omega)
\]
and summing the resulting geometric series yields
\[
S(\omega)
=
\frac{2\bigl(32\cos^3\omega+12\cos^2\omega-116\cos\omega-105\bigr)}
{9\,(4\cos\omega-5)\,(4\cos\omega+5)^2}.
\]
At \(\omega=0\) this gives \(S(0)=118/243\), in agreement with the CLT variance.


% Cut block: Partition-function recurrence and discriminant (Appendix D)
\section{The partition-function recurrence and discriminant}
\label{app:recurrence}

This appendix collects the transfer-matrix calculations used in Section~\ref{sec:partition}. Once the four-state presentation from Appendix~\ref{app:pressure} is fixed, the remaining arguments are linear algebra and symbolic elimination, so we separate them from the main narrative.

\subsection{Weighted transfer matrix and boundary vectors}
Weight each edge of the four-state graph from Appendix~\ref{app:pressure} by \(y^{Y}\), where \(Y\in\{0,1\}\) is the second component of the label \((X,Y)\). In the state order \((00,01,10,11)\) this gives the weighted adjacency matrix
\[
M(y)=
\begin{pmatrix}
1&y&0&y\\
0&0&1&0\\
1&1+y&0&0\\
1&0&0&0
\end{pmatrix}.
\]
Its characteristic polynomial is
\[
\det(\lambda I-M(y))
=
\lambda^4-\lambda^3-(2y+1)\lambda^2+\lambda+y(y+1)
=
\Pi(\lambda,y).
\]
Replacing \(\lambda\) by \(t^{-1}\) gives
\[
\det(I-tM(y))
=
1-t-(2y+1)t^2+t^3+y(y+1)t^4
=
D(t;y).
\]
Because the synchronized graph is finite, the weighted language counted by \(Z_m(y)\) is rational and admits a four-dimensional linear representation; compare the standard transfer-matrix formalism in \cite[Chapter~6]{AlloucheShallit2003} and \cite{GrabnerKirschenhoferTichy2002}. With \(\alpha\) fixed by the synchronized initial state \(00\), a convenient terminal vector is obtained by solving
\[
\alpha^{\mathsf T}M(y)^m\beta(y)=Z_m(y)\qquad (m=0,1,2,3),
\]
using the directly enumerated initial polynomials
\[
Z_0(y)=1,\qquad
Z_1(y)=1+y,\qquad
Z_2(y)=2+2y,\qquad
Z_3(y)=2+5y+y^2.
\]
The resulting representation is
\[
Z_m(y)=\alpha^{\mathsf T}M(y)^m\beta(y),
\qquad
\alpha=
\begin{pmatrix}
1\\ 0\\ 0\\ 0
\end{pmatrix},
\qquad
\beta(y)=
\begin{pmatrix}
1\\[1mm]
\dfrac{1}{1+y}\\[2mm]
\dfrac{1}{y}\\[2mm]
\dfrac{y}{1+y}
\end{pmatrix}.
\]
The rational entries of \(\beta(y)\) are an artifact of this compact \(4\times4\) representation; after multiplication by \(M(y)^m\) all poles cancel and one recovers the polynomial \(Z_m(y)\).

\subsection{Numerator computation}
The generating function \(Z(t;y)=\sum_{m\ge 0}Z_m(y)t^m\) satisfies
\[
Z(t;y)=\frac{N(t;y)}{D(t;y)},
\]
where \(N(t;y)=Z(t;y)\cdot D(t;y)\) is determined by the initial values. Setting
\[
\begin{aligned}
N(t;y)
&=\sum_{m=0}^{3}
\Bigl(
Z_m(y)-\sum_{k=1}^{\min(m,4)}c_k(y)Z_{m-k}(y)
\Bigr)t^m,
\end{aligned}
\]
where \(c_1=1\), \(c_2=2y+1\), \(c_3=-1\), \(c_4=-y(y+1)\) are the recurrence coefficients, a direct calculation using \(Z_0=1\), \(Z_1=1+y\), \(Z_2=2+2y\), \(Z_3=2+5y+y^2\) gives
\[
N(t;y)=1+yt-yt^2-y^2t^3=(1+ty)(1-yt^2).
\]
The factorization is verified by expanding the product.

\subsection{Recurrence and Hankel rank}
Because \(Z_m(y)\) is obtained by summing the weights \(y^{\abs{\Fold_m(\omega)}_1}\) over all length-\(m\) microstate words, it is a linear functional of the weighted transfer-matrix dynamics. Consequently \(Z_m(y)\) satisfies the characteristic recurrence of \(M(y)\),
\[
Z_m(y)=Z_{m-1}(y)+(2y+1)Z_{m-2}(y)-Z_{m-3}(y)-y(y+1)Z_{m-4}(y)
\qquad (m\ge 4).
\]
The initial values
\[
Z_0(y)=1,\qquad
Z_1(y)=1+y,\qquad
Z_2(y)=2+2y,\qquad
Z_3(y)=2+5y+y^2
\]
come from direct enumeration. Therefore
\[
Z(t;y)=\sum_{m\ge0}Z_m(y)t^m
\]
is a rational function with denominator \(D(t;y)\) and numerator \(N(t;y)=(1+ty)(1-yt^2)\).

\subsection{Discriminant factorization}
Let
\[
\Pi(\lambda,y)=\lambda^4-\lambda^3-(2y+1)\lambda^2+\lambda+y(y+1).
\]
A direct symbolic computation gives
\[
\Disc_\lambda \Pi(\lambda,y)
=-y(y-1)\bigl(256y^3+411y^2+165y+32\bigr).
\]
The cubic factor has exactly one negative real root,
\[
y_{\mathrm{c}}\approx -1.13445.
\]
This is the unique negative real value of \(y\) at which \(\Pi(\lambda,y)\) has a repeated root, apart from the zeros at \(y=0\) and \(y=1\).

\subsection{Root analysis at the discriminant zeros}
At each discriminant zero, we compute the roots of \(\Pi(\lambda,y)\) explicitly.

\emph{At \(y=0\):} \(\Pi(\lambda,0)=\lambda^4-\lambda^3-\lambda^2+\lambda=\lambda(\lambda-1)^2(\lambda+1)\). The repeated root is \(\lambda=1\), which is dominant. Since \(Z(t;0)=1/((1-t)^2(1+t))\), the coefficients \(Z_m(0)=\abs{\Fold_m^{-1}(0^m)}\) grow linearly in \(m\), exactly as expected from the double dominant root.

\emph{At \(y=1\):} \(\Pi(\lambda,1)=\lambda^4-\lambda^3-3\lambda^2+\lambda+2=(\lambda-2)(\lambda+1)^2(\lambda-1)\). The repeated root is \(\lambda=-1\), which is subdominant (dominated by \(\lambda=2\)).

\emph{At \(y=y_{\mathrm{c}}\approx -1.13445\):} the quartic \(\Pi(\lambda,y_{\mathrm{c}})\) evaluates to approximately \(\lambda^4-\lambda^3+1.2689\lambda^2+\lambda+0.1525\). Numerical root-finding gives a repeated real root near \(\lambda\approx -0.2734\) and a complex-conjugate pair near \(\lambda\approx 0.7734\pm 1.2008\,i\) with modulus approximately \(1.428\). Since \(|-0.2734|<1.428\), the repeated root is strictly subdominant. Hence the discriminant zero at \(y_{\mathrm{c}}\) does not produce a singularity of the dominant eigenvalue.

\begin{remark}
For \(y>0\), the Perron--Frobenius theorem guarantees that the dominant root of \(M(y)\) is simple, real, and positive. The free energy \(f(y)=\log\lambda_\star(y)\) is therefore real-analytic on \((0,\infty)\). On the negative real axis, \(Z_m(y)\) involves sign-alternating terms and a separate dominant-sheet analysis would be needed; we do not pursue this.
\end{remark}

\subsection{Output-weight cumulants}
Set \(y=e^\theta\). The scaled cumulant generating function for \(H_m\) is
\[
\Lambda_Y(\theta)
=
\log \rho\!\bigl(M(e^\theta)\bigr)-\log 2.
\]
Write
\[
F(\lambda,z):=\lambda^4-\lambda^3-(2z+1)\lambda^2+\lambda+z(z+1).
\]
Since \(F(\rho(\theta),e^\theta)=0\) and \(F(2,1)=0\), implicit differentiation at \(\theta=0\) gives
\[
F_\lambda(2,1)\rho'(0)+F_z(2,1)=0.
\]
Here
\[
F_\lambda(2,1)=9,\qquad F_z(2,1)=-5,
\]
so \(\rho'(0)=5/9\). A second differentiation yields
\[
F_{\lambda\lambda}(2,1)\rho'(0)^2+2F_{\lambda z}(2,1)\rho'(0)+F_{zz}(2,1)+F_z(2,1)+F_\lambda(2,1)\rho''(0)=0.
\]
Using
\[
F_{\lambda\lambda}(2,1)=30,\qquad
F_{\lambda z}(2,1)=-8,\qquad
F_{zz}(2,1)=2,
\]
one obtains \(\rho''(0)=71/243\). Consequently,
\[
\Lambda_Y'(0)=\frac{5}{18},
\qquad
\Lambda_Y''(0)=\frac{67}{972}.
\]
These are the mean density and asymptotic variance appearing in Corollary~\ref{thm:weight-stats}. Since \(\theta\mapsto \rho(M(e^\theta))\) is real-analytic on \(\RR\), the G\"artner--Ellis rate function is
\[
I(a)=\sup_{\theta\in\RR}\bigl(\theta a-\Lambda_Y(\theta)\bigr).
\]
